\documentclass[11pt,a4paper]{article}

\usepackage[T1]{fontenc}
\usepackage[utf8]{inputenc}
\usepackage{charter}
\usepackage[scaled=0.94]{helvet}
\usepackage[left=1.05in,right=1.05in,top=0.95in,bottom=0.95in]{geometry}
\usepackage{parskip}
\usepackage[hidelinks]{hyperref}
\usepackage{microtype}

\renewcommand{\familydefault}{\rmdefault}
\linespread{1.05}
\pagestyle{empty}

\newcommand{\sender}[1]{{\small\sffamily #1\par}}

\begin{document}

% ---------------------------------------------------------------- sender block
{\sffamily\bfseries\large Parth Kumar Sinha}\par
\vspace{2pt}
\sender{Independent Researcher}
\sender{Oxford, United Kingdom}
\sender{\href{mailto:sinhaparth555@gmail.com}{sinhaparth555@gmail.com}}
\sender{ORCID \href{https://orcid.org/0009-0002-3120-9301}{0009-0002-3120-9301}}

\vspace{18pt}
\hrule
\vspace{18pt}

14 September 2026

\vspace{10pt}

The Editor-in-Chief\\
\textit{IEEE Transactions on Parallel and Distributed Systems}

\vspace{14pt}

\textbf{Re:} Submission of ``Portable, Seekable Random Streams for Parallel CPU
Simulation'' as a regular paper

\vspace{10pt}

Dear Editor-in-Chief,

I am submitting ``Portable, Seekable Random Streams for Parallel CPU
Simulation'' for consideration as a regular paper in IEEE Transactions on
Parallel and Distributed Systems.

The paper starts from a defect rather than a design. Checkpointing a
\texttt{std::mt19937} and restoring it under a different C++ standard library is
not reliable: the libraries disagree on the numeric base of the serialized
digits and on the order of the 624 state words. Under the Microsoft standard
library the restore fails outright. Under libc++ it succeeds and resumes at the
wrong position, with no diagnostic, which is the case that worries me. A user
report against XGBoost 3.3.0 shows the first failure happening in the field.

Storing a position instead of an implementation removes the problem by
construction, and for a counter-based generator the position is a key and an
integer. What has kept that substitution out of CPU codebases is a performance
objection. My own proposal to make the serialized generator state of XGBoost
portable was declined on the grounds that Philox runs at one tenth the speed of
the Mersenne Twister. The paper quotes that review comment verbatim and
discloses in a footnote that I submitted the proposal it declined. The
evaluation exists to answer that sentence.

It does not reproduce. Across five processors, unspecialised scalar Philox costs
between 1.10 and 1.65 times what \texttt{std::mt19937} costs per byte, not ten
times. With the wide multiplies specialised and counters interleaved across SIMD
lanes, the bulk path costs between 0.22 and 0.75 times \texttt{std::mt19937}. In
parallel, aggregate cost per byte does not move with the thread count, and the
first limit on a sixteen-core host is hyperthread co-location rather than memory
bandwidth.

Three things about the evaluation may interest the reviewers. Instruction supply
is excluded by measurement rather than by argument: co-location costs 59 percent
more per byte while L1 instruction-cache misses fall by 36 percent. The flat
scaling result is reproduced under a second, independent threading runtime, so
it is a property of the generator and not of my worker pool. And every placement
study runs behind a gate that first measures whether the host can place threads
at all, because a virtual machine accepts every affinity mask while the
hypervisor places vCPUs itself, which silently turns two arms into one
measurement taken twice.

I have tried to be straightforward about what the measurements do not show.
xoshiro256++ is faster than vphilox on four of the five machines, and on ARM the
buffered engine interface is still slower than \texttt{std::mt19937} even though
the NEON kernel beneath it is faster. Both are in the paper. Execution-port
contention is reported as an inference from the instruction mix rather than as a
measurement, since the counters that would test it directly are not available on
the one host that can control placement.

Every number traces to an archived run. The source, the raw JSON and the script
that regenerates every table and figure are at
\url{https://github.com/sinhaparth5/vphilox}, the measured release is archived
at \href{https://doi.org/10.5281/zenodo.22103483}{doi:10.5281/zenodo.22103483},
and continuous integration fails the build if a derived file no longer matches
the data behind it.

The manuscript is original, is not under consideration elsewhere, and has not
been published in whole or in part. I am an independent researcher with no
funding to declare. Beyond the XGBoost disclosure above I have no competing
interests.

Thank you for your time.

\vspace{16pt}

Yours sincerely,

\vspace{22pt}

Parth Kumar Sinha

\end{document}
